\documentclass[10pt,a4paper]{amsart}
\usepackage[margin=19mm]{geometry}
\usepackage{amsmath,amssymb,mathtools}
\usepackage[scr=rsfs,cal=euler]{mathalfa}
\usepackage{xfrac}
\usepackage[unicode,hidelinks,bookmarks=false]{hyperref}
\newcommand{\ie}{i.e.\ }
\newcommand{\cf}{cf.\ }
\newcommand{\resp}{respectively\ }
\newcommand{\Subsection}{\subsection}
\providecommand{\nobreakdash}{\nobreak\hbox{-}}
\setlength{\emergencystretch}{2em}
\multlinegap0pt
\usepackage{xcolor}
\usepackage{amssymb}
\usepackage{mathtools}
\newtheorem{lemma}{Lemma}
\newcommand{\PP}{\R[z]_{<m}}
\newcommand{\PQ}{\Q[z]_{<m}}
\newcommand{\Q}{\mathbb{Q}}
\newcommand{\R}{\mathbb{R}}
\newcommand{\defas}{\coloneqq}
\title{Values of Absorbing Recursive Games Express\\ All Real Algebraic Numbers}
\author{Ali Asadi, Krishnendu Chatterjee}
\date{Source: arXiv:2609.11583v2 (CC BY 4.0)}
\begin{document}
\maketitle

\noindent\emph{Source and licence.} Values of Absorbing Recursive Games Express\\ All Real Algebraic Numbers. Version arXiv:2609.11583v2 (CC BY 4.0), distributed by arXiv under the Creative Commons Attribution 4.0 International licence, http://creativecommons.org/licenses/by/4.0/, as recorded in the arXiv licence metadata for this exact version and shown on the abstract page https://arxiv.org/abs/2609.11583v2. Full licence notice: \texttt{licenses/arxiv-2609.11583-CC-BY-4.0.txt}.\par\smallskip

\noindent\textit{Editorial scope.} Counted unit: the article's lemma lem:basis (source0.tex lines 497--504). The article's complete original proof of it is delivered with it (source0.tex lines 506--566, one \texttt{proof} environment of 61 lines). No mathematics is written, repaired or re-derived: the statement and the proof are the article's own lines, in the article's own notation, and no retained line is substituted or re-flowed.  No further source-local context is needed: every cross-reference of every retained line resolves inside the two delivered spans.  Nothing was omitted from any retained span, so the package carries no omission record.  No retained line contains a citation, so the wrapper carries no bibliography and no bibliography entry.  No retained line contains a figure environment, an image-inclusion command or drawing code, so no image is delivered and the package declares no asset dependency.  Editorial formatting only, in addition: an AMS \texttt{amsart} preamble that restates the article's own declarations and macros used by the retained lines, namely the environments \texttt{lemma} on separate counters, together with the standard packages those lines need.  The retained environments print in this document as Lemma 1 in order of appearance, whereas the article prints the same items with the numbering of its own full document.  The complete native source of the article as distributed in the arXiv e-print for this exact version is bundled unchanged as \texttt{sources/210034/source0.tex}.
\begin{lemma}\label{lem:basis}
There exist $p_1,\dots,p_m\in\PQ$ and real numbers $\tilde y_1,\dots,\tilde y_m$ such that:
\begin{enumerate}
\item[\emph{(i)}] $(p_1,\dots,p_m)$ is a basis of $\PP$ over $\R$ (hence of $\PQ$ over $\Q$);
\item[\emph{(ii)}] $p_k(\alpha)>0$ for every $k$;
\item[\emph{(iii)}] $\ell=\sum_{k=1}^m\tilde y_kp_k$ with $\tilde y_k>0$ for every $k$.
\end{enumerate}
\end{lemma}

\begin{proof}
Consider the set
\begin{align*}
\mathcal U\defas \Bigl\{(f_1,\dots,f_m)\in(\PP)^m:\ &(f_k)_k\text{ is a basis of }\PP\\
&\land f_k(\alpha)>0\ \forall k\\
&\land \text{the coordinates of }\ell\text{ in }(f_k)_k\text{ are all }>0\Bigr\}.
\end{align*}

We first show that $\mathcal U$ is open in $(\PP)^m$.
Indeed, let $(f_1,\dots,f_m)$ be any tuple in $(\PP)^m$ and let $F\in\R^{m\times m}$ be the matrix whose $k$-th column is the coefficient vector of $f_k$.
Note that $\mathcal U$ is constrained by the three conditions in its definition.
We check that each is an open condition on $F$, so that $\mathcal U$, their intersection, is open.
First, $(f_1,\dots,f_m)$ is a basis if and only if $\det F\neq0$, which is open because $\det$ is continuous.
Second, each value $f_k(\alpha)$ is a linear, hence continuous, function of $F$, so the requirement $f_k(\alpha)>0$ for all $k$ is open.
Third, on the open set where $F$ is invertible, the coordinates $c_1,\dots,c_m$ of the fixed polynomial $\ell$, defined by $\ell=\sum_kc_kf_k$, are given by Cramer's rule as continuous functions of $F$, so the requirement $c_k>0$ for all $k$ is open.
Therefore, $\mathcal U$ is open in $(\PP)^m$.

We now show that $\mathcal U$ is nonempty. Let
\[
\mathcal Z\defas \{f\in\PP:f(\alpha)=0\}.
\]
This is the kernel of the nonzero linear map $f\mapsto f(\alpha)$. Hence, we have
$\dim\mathcal Z=m-1$. Also $\ell\notin\mathcal Z$, since $\ell(\alpha)=1$.
Choose a basis $h_1,\dots,h_{m-1}$ of $\mathcal Z$ (an empty list if $m=1$), and set
\[
u_k\defas h_k\ (1\leq k\leq m-1),
\qquad
u_m\defas -(h_1+\dots+h_{m-1}),
\]
with the convention that the empty sum is $0$. Then $u_1,\dots,u_m\in\mathcal Z$ and
\[
u_1+\dots+u_m=0.
\]
Define $f_k\defas \ell+u_k$ for $1\leq k\leq m$. Since $u_k(\alpha) = 0$, we have
\[
f_k(\alpha)=\ell(\alpha)+u_k(\alpha)=1>0
\qquad(1\leq k\leq m).
\]
Moreover, we have
\begin{align*}
    \frac1m\sum_{k=1}^mf_k &= \frac1m\sum_{k=1}^m(\ell+u_k)\\
    &= \frac1m\sum_{k=1}^m\ell + \frac1m\sum_{k=1}^mu_k\\
    &= \ell .
\end{align*}
Thus, if $(f_1,\dots,f_m)$ is a basis, then the coordinates of $\ell$ in this basis are all equal to $\frac1m>0$.

It remains to prove that $(f_1,\dots,f_m)$ is a basis. Suppose
$\sum_{k=1}^mc_kf_k=0$. Evaluating at $\alpha$ gives $\sum_kc_k=0$. Hence
\[
0=\sum_{k=1}^mc_kf_k
=\Bigl(\sum_{k=1}^mc_k\Bigr)\ell+\sum_{k=1}^mc_ku_k
=\sum_{k=1}^mc_ku_k
=\sum_{k=1}^{m-1}(c_k-c_m)h_k .
\]
Since $h_1,\dots,h_{m-1}$ are linearly independent, $c_k=c_m$ for all
$k\leq m-1$. Together with $\sum_kc_k=0$, this gives $c_1=\dots=c_m=0$.
Therefore $(f_1,\dots,f_m)$ is a basis.
Hence $(f_1,\dots,f_m)\in\mathcal U$.

Since $(\PQ)^m$ is dense in $(\PP)^m$, the nonempty open set $\mathcal U$ contains a rational tuple $(p_1,\dots,p_m)\in(\PQ)^m$. Properties (ii) and (iii) hold by definition of $\mathcal U$. The tuple $(p_1,\dots,p_m)$ is a basis of $\PP$ over $\R$ by definition of $\mathcal U$; since the $p_k$ are rational polynomials and are linearly independent over $\R$, they are linearly independent over $\Q$, and as $\dim_\Q\PQ=m$, they form a basis of $\PQ$ over $\Q$. Thus property (i) holds, which completes the proof.
\end{proof}

\end{document}
